\documentclass{kazetex}

\title{高等数学习题集}
\author{风待桐叶}
\date{2026年5月}

\begin{document}

\maketitle

\tableofcontents
\newpage

\section{极限与连续}

\subsection{极限的计算}

\begin{problem}
  求极限 $\displaystyle\lim_{x\to 0}\frac{\sin x}{x}$。
\end{problem}

\begin{problem}
  求极限 $\displaystyle\lim_{x\to\infty}\left(1+\frac{1}{x}\right)^x$。
\end{problem}

\begin{problem}
  求极限 $\displaystyle\lim_{x\to 0}\frac{\tan x - \sin x}{x^3}$。
\end{problem}

\subsection{重要极限}

\begin{problem}
  求极限 $\displaystyle\lim_{x\to 0}\frac{1-\cos x}{x^2}$。
\end{problem}

\begin{problem}
  求极限 $\displaystyle\lim_{x\to 0}(1+2x)^{\frac{1}{x}}$。
\end{problem}

\subsection{连续性}

\begin{problem}
  判断函数 $f(x)=\dfrac{|x|}{x}$ 在 $x=0$ 处是否连续。
\end{problem}

\begin{problem}
  设 $f(x)=\begin{cases}
    x^2\sin\frac{1}{x}, & x\neq 0\\
    0, & x=0
  \end{cases}$，讨论 $f(x)$ 在 $x=0$ 处的连续性与可导性。
\end{problem}

\section{导数与微分}

\subsection{基本求导法则}

\begin{problem}
  求 $y=x^3-3x^2+2x-1$ 的导数。
\end{problem}

\begin{problem}
  求 $y=\ln(\sin x)$ 的导数。
\end{problem}

\subsection{高阶导数}

\begin{problem}
  设 $y=\mathrm{e}^{x}\cos x$，求 $y''$。
\end{problem}

\begin{problem}
  设 $y=\dfrac{1}{1-x}$，求 $y^{(n)}$。
\end{problem}

\subsection{导数的应用}

\begin{problem}
  求曲线 $y=x^3-3x$ 的拐点与凹凸区间。
\end{problem}

\begin{problem}
  求函数 $f(x)=x^3-3x^2-9x+5$ 在 $[-2,4]$ 上的最大值与最小值。
\end{problem}

\subsection{微分中值定理}

\begin{problem}
  设 $f(x)$ 在 $[0,1]$ 上连续，$(0,1)$ 内可导，且 $f(0)=0$，$f(1)=1$，
  证明：存在 $\xi\in(0,1)$ 使得 $f'(\xi)=2\xi$。
\end{problem}

\begin{problem}
  验证函数 $f(x)=\ln x$ 在 $[1,\mathrm{e}]$ 上满足 Lagrange 中值定理的条件，并求 $\xi$。
\end{problem}

\section{不定积分}

\subsection{分部积分法}

\begin{problem}
  计算 $\displaystyle\int x\,\mathrm{e}^{x}\,\mathrm{d}x$。
\end{problem}

\begin{problem}
  计算 $\displaystyle\int \ln x\,\mathrm{d}x$。
\end{problem}

\subsection{有理函数积分}

\begin{problem}
  计算 $\displaystyle\int \frac{1}{x^2-1}\,\mathrm{d}x$。
\end{problem}

\begin{problem}
  计算 $\displaystyle\int \frac{x+1}{x^2+4x+3}\,\mathrm{d}x$。
\end{problem}

\subsection{三角函数的积分}

\begin{problem}
  计算 $\displaystyle\int \sin^2 x\,\mathrm{d}x$。
\end{problem}

\begin{problem}
  计算 $\displaystyle\int \frac{1}{1+\cos x}\,\mathrm{d}x$。
\end{problem}

\subsection{换元法}

\begin{problem}
  计算 $\displaystyle\int \frac{\ln x}{x}\,\mathrm{d}x$。
\end{problem}

\begin{problem}
  计算 $\displaystyle\int \frac{1}{\sqrt{1-x^2}}\,\mathrm{d}x$。
\end{problem}

\end{document}
